% ==============================================================================
% T0/FFGFT Theory: Document 232
% Title: Quantum Graphity as a subset of FFGFT
% Author: Johann Pascher
% Date: May 2026
% References: Documents 230 (Hilbert-space translation), 231 (Hilbert-space
%             extensions), 202 (field theory), 210 (windings).
%             Source paper: Konopka, Markopoulou, Severini,
%             Phys. Rev. D 77, 104029 (2008), arXiv:0801.0861
% ==============================================================================

\section*{Preamble}

Documents~230 and~231 showed that standard quantum mechanics
(Hilbert-space formalism) is a \emph{subset} of FFGFT, arising from
the full geometry through a reduction chain
\[
S^2 \subset \text{cylinder} \subset T^2 \subset T^3 \subset T^4.
\]
This document applies the same methodological strategy to another
established theory: \textbf{Quantum Graphity} (Konopka, Markopoulou,
Severini, 2008).

Quantum Graphity is a peer-reviewed, background-independent model
in which spacetime and geometry arise as emergent properties of a
fundamental \emph{graph}. It is mathematically clearly defined,
published in \emph{Phys.~Rev.~D}~77, 104029 (2008), and thus a
concrete object of comparison.

\textbf{Central thesis of this document:} Quantum Graphity \emph{can
plausibly be sketched as a subset of FFGFT}, through a five-stage
reduction chain from the full FFGFT geometry -- analogous to the
Bloch sphere, but through different reduction steps. Unlike
Doc.~230, however, this translation is \emph{not carried out} but a
hypothesis: the mathematical triangulation construction and the
associated convergence proof are open research tasks (see §10.4
below). What is methodologically important already is that FFGFT
shares with Quantum Graphity (and LQG, CDT, Wolfram hypergraphs) the
basic programme \enquote{emergent geometry from discrete
microstructure} -- without inheriting their limitations.

% ==============================================================================
\section*{Notation}
% ==============================================================================

\subsection*{Quantum Graphity symbols}

\begin{center}
\renewcommand{\arraystretch}{1.3}
{\small
\adjustbox{max width=\textwidth}{%
\begin{tabular}{>{\raggedright\arraybackslash}p{3.6cm}>{\raggedright\arraybackslash}p{7.4cm}}
\toprule
\textbf{Symbol} & \textbf{Meaning} \\
\midrule
$K_N$ & complete graph on $N$ vertices (all vertices connected) \\
$N$ & number of vertices (typically $N \gg 1$) \\
$V_a$ & vertex (node) $a$ \\
$e_{ab}$ & edge between vertex $a$ and $b$ \\
$n_{ab} \in \{0, 1\}$ & occupation number of edge $e_{ab}$ (off/on), base model \\
$|0\rangle, |1\rangle$ & edge states: off/on \\
$a^\dagger_{ab}, a_{ab}$ & creation and annihilation operators on edge $e_{ab}$ \\
$\{a, a^\dagger\}$ & fermionic anticommutator, $= 1$ \\
$N_{ab} = a^\dagger_{ab} a_{ab}$ & number operator (adjacency matrix element) \\
$|j, m\rangle$ & extended edge labels ($j = 0, 1$, $m \in \{-j, \ldots, j\}$) \\
$M, M^\pm$ & angular momentum operators on $|j, m\rangle$ labels \\
$\mathcal{H}_{\text{edge}}$ & Hilbert space of a single edge \\
$\mathcal{H}_{\text{total}}$ & tensor product Hilbert space of all edges \\
$H_V, H_B, H_C, H_D, H_\pm$ & Hamiltonian components (valence, cycles, $C$, $D$, loop) \\
$g_V, g_B, g_C, g_D, g_\pm$ & coupling constants of Hamiltonian terms \\
$v_0$ & preferred vertex valence (typically $v_0 = 3$ for hexagonal) \\
$S_N$ & permutation group on $N$ elements \\
\bottomrule
\end{tabular}%
}
}
\renewcommand{\arraystretch}{1.0}
\end{center}

\subsection*{FFGFT symbols (recalled from Doc.~230/231)}

\begin{center}
\renewcommand{\arraystretch}{1.3}
{\small
\adjustbox{max width=\textwidth}{%
\begin{tabular}{>{\raggedright\arraybackslash}p{3.6cm}>{\raggedright\arraybackslash}p{7.4cm}}
\toprule
\textbf{Symbol} & \textbf{Meaning} \\
\midrule
$T^n$ & $n$-torus, $n$ compact loops \\
$T^4$ & full FFGFT geometry incl.\ time winding \\
$(n_x, n_y, n_z, k_t)$ & FFGFT winding quantum numbers \\
$\varphi_{n,l,j}$ & FFGFT mode function on $T^4$ \\
$\xi_0 = 4/30000$ & FFGFT geometry parameter \\
$L_0 = \xi_0 \cdot L_P$ & FFGFT minimum length scale (sub-Planck) \\
$L_P$ & Planck length \\
$U(1)$ & electromagnetic gauge group \\
\bottomrule
\end{tabular}%
}
}
\renewcommand{\arraystretch}{1.0}
\end{center}

\subsection*{Other symbols}

\begin{center}
\renewcommand{\arraystretch}{1.3}
{\small
\adjustbox{max width=\textwidth}{%
\begin{tabular}{>{\raggedright\arraybackslash}p{3.6cm}>{\raggedright\arraybackslash}p{7.4cm}}
\toprule
\textbf{Symbol} & \textbf{Meaning} \\
\midrule
$\subset$ & subset of, contained in \\
$\to$ & reduction or transition \\
$\beta = 1/k_BT$ & inverse temperature \\
$T_c$ & critical temperature (phase transition) \\
\bottomrule
\end{tabular}%
}
}
\renewcommand{\arraystretch}{1.0}
\end{center}

% ==============================================================================
\section{Quantum Graphity in one page}
% ==============================================================================

Quantum Graphity (Konopka, Markopoulou, Severini, \emph{Phys.~Rev.~D}~77,
104029 (2008), arXiv:0801.0861) is a background-independent model in
which space and geometry emerge from a dynamical \emph{graph}. The
basic assumptions are:

\begin{itemize}[leftmargin=2em,itemsep=2pt,topsep=2pt]
  \item \textbf{Complete graph $K_N$:} on $N \gg 1$ vertices, \emph{all}
        vertex pairs are connected by edges. $K_N$ has $\binom{N}{2}$
        edges.
  \item \textbf{Edge Hilbert space (base model):} each edge $e_{ab}$
        carries a Hilbert space
        $\mathcal{H}_{\text{edge}} = \text{span}\{|0\rangle, |1\rangle\}$
        -- a \emph{fermionic oscillator} (or equivalently a qubit).
        $|0\rangle$ is \enquote{off}, $|1\rangle$ is \enquote{on}.
  \item \textbf{Tensor-product Hilbert space:} the total Hilbert
        space is
        $\mathcal{H}_{\text{total}} = \bigotimes_{a<b} \mathcal{H}_{\text{edge}}$,
        i.e.\ tensor product over all $\binom{N}{2}$ edges.
  \item \textbf{Adjacency operators:} $N_{ab} = a^\dagger_{ab} a_{ab}$
        is the quantum-mechanical adjacency matrix; the active edges
        form a subgraph of $K_N$.
  \item \textbf{Hamiltonian:} several permutation-invariant terms
        ($H_V$ valence, $H_B$ cycles, etc.) drive a \emph{phase
        transition}.
  \item \textbf{High-energy phase:} all edges on, full permutation
        symmetry $S_N$, no notion of locality.
  \item \textbf{Low-temperature phase:} permutation symmetry breaks
        to translation symmetry. A \emph{regular subgraph} emerges
        -- in the standard case a \emph{hexagonal}
        ($3$-regular) configuration.
  \item \textbf{Extension with \enquote{matter labels}
        $|j, m\rangle$:} in §IV of the paper the base model is
        extended: each edge can carry $j = 0, 1$ and
        $m \in \{-j, \ldots, +j\}$. This extension allows
        \emph{emergent matter} and \emph{emergent $U(1)$ gauge theory}
        through string-net condensation (Levin-Wen mechanism).
\end{itemize}

\textbf{What Quantum Graphity does not do:} it does not deliver
Standard-Model constants, fermion masses, or the fine-structure
constant. It delivers the skeleton \enquote{how geometry emerges
from discrete microstructure}, plus an emergent $U(1)$ gauge
theory, without specifically explaining the SM matter structure.

% ==============================================================================
\section{Methodological parallel to Hilbert-space reduction}
% ==============================================================================

In Doc.~230 we showed how the standard Bloch sphere of quantum
mechanics arises from FFGFT through a geometric reduction
($T^4 \to T^3 \to T^2 \to \text{cylinder} \to S^2$). At each step a
specific structural layer was lost -- fractal corrections, winding
quantum numbers, bundle structure, time winding.

Here we exhibit the analogous reduction for Quantum Graphity:
\begin{equation}
  \boxed{\text{Quantum Graphity} \subset \text{FFGFT}}
\end{equation}
through a different -- but structurally comparable -- reduction
chain. The reduction steps are:
\begin{enumerate}[leftmargin=2em,itemsep=2pt,topsep=2pt]
  \item \textbf{Discretisation:} continuous $T^4$ $\to$ discrete
        graph with $N$ vertices.
  \item \textbf{Topology forgetting:} $T^4$ topology $\to$ abstract
        complete graph $K_N$.
  \item \textbf{Quantum-number generalisation:} FFGFT windings
        $(n_x, n_y, n_z, k_t)$ $\to$ general edge occupation /
        spin labels.
  \item \textbf{Parameter generalisation:} $\xi_0 = 4/30000$ $\to$
        general coupling parameters.
  \item \textbf{Matter dropping:} 12 fermion masses + $\alpha$
        $\to$ no counterpart.
\end{enumerate}

These five reductions together yield the Quantum-Graphity skeleton.
They do \emph{not} correspond to the Bloch-sphere reduction of
Doc.~230 -- they are a \emph{different} projection of the same
FFGFT full geometry onto a different abstract description.

% ==============================================================================
\section{Reduction 1: continuum to discrete}
% ==============================================================================

\subsection{\texorpdfstring{FFGFT: continuous windings on $T^4$}{FFGFT: continuous windings on T4}}

In FFGFT the full geometry lives on a 4D-torus $T^4 = T^3 \times S^1$.
Mode functions $\varphi_{n,l,j}(x_0, x_1, x_2, t)$ are \emph{smooth
functions} on this compact manifold. The mode quantum numbers
$(n_x, n_y, n_z, k_t)$ are discrete -- but each mode is a continuous
function, and the underlying geometry is differential-geometrically
smooth.

\subsection{Quantum Graphity: discrete graph}

In Quantum Graphity there is no smooth manifold. There is only the
discrete complete graph $K_N$ with $N$ vertices and $\binom{N}{2}$
possible edges. Each edge is either active or inactive; active
edges may carry labels.

\subsection{Reduction}

Discretisation \emph{would} be sketched formally as triangulation of
$T^4$: subdividing $T^4$ into $N$ small cells (e.g.\ tetrahedra or
small cubes) yields a graph whose vertices correspond to cells and
whose edges encode neighbourhood relations. Heuristically, as the
mesh becomes fine ($N \to \infty$, cell size $\to 0$), this graph
should converge to the continuous manifold.

\textbf{Important caveat:} this construction is a \emph{sketch}, not
a carried-out proof. A rigorous triangulation translation FFGFT
$\to$ Quantum Graphity has not been worked out in the FFGFT
documents -- it is an open research task (see §10.4 below). The
methodological plausibility that such a construction is possible
rests on the established tradition of simplicial approximation in
algebraic topology. A concrete convergence theorem for the FFGFT
mode functions $\varphi_{n,l,j}$ under a specific $T^4$ triangulation,
however, is not proved.

\textbf{What is lost:} differential geometry. The FFGFT equation of
motion $\mathcal{D}_{T^4}\varphi = E\varphi$ (Dirac-like operator on
$T^4$) has no direct counterpart in the discrete picture -- one
obtains instead a Hamiltonian update rule on the edges.

\textbf{What remains:} the topological information \emph{if} one
remembers that the graph came from a 4D-torus triangulation. In
Quantum Graphity, however, this memory is \emph{deliberately
forgotten} (see Reduction~2).

% ==============================================================================
\section{Reduction 2: forgetting topology}
% ==============================================================================

\subsection{\texorpdfstring{FFGFT: $T^4$ topology}{FFGFT: T4 topology}}

The 4D-torus has a very specific topology: four compact loops, each
with winding number $\in \mathbb{Z}$. This topology is \emph{part of
the physical statement} of FFGFT -- the winding quantum numbers are
not arbitrary but correspond to specific physical quantities
(masses, charges, spin).

\subsection{\texorpdfstring{Quantum Graphity: complete graph $K_N$}{Quantum Graphity: complete graph KN}}

In Quantum Graphity one starts with the \emph{complete} graph $K_N$
-- thus with \emph{maximal} connectivity, without any topology. This
is the high-temperature state. Only the phase transition at low
temperature lets a low-dimensional topology emerge.

\textbf{From the FFGFT perspective:} the complete graph $K_N$
corresponds to \enquote{forgetting} the $T^4$ topology. If one no
longer remembers that the $N$ points were originally organised in a
4D-torus geometry, every point can be connected to every other --
yielding $K_N$.

\subsection{Reduction}

Heuristically the step can be described as: starting from a
(sketched) triangulation of $T^4$ from Reduction~1, one would have a
\emph{local} graph $G_{T^4}$ in which each vertex has only nearby
neighbours. Stripping this locality information by declaring all
vertex pairs as potentially connectable would yield $K_N$.

\textbf{This step too remains a sketch.} The concrete construction
of a map $G_{T^4} \to K_N$ and the proof that under this map the
Quantum-Graphity Hamiltonian components ($H_V$, $H_B$, etc.)
reproduce the correct FFGFT low-energy physics are not worked out.

\textbf{What is lost:} locality. In FFGFT, locality is given (by
the $T^4$ geometry). In Quantum Graphity, locality is an
\emph{emergent phenomenon} of the low-temperature phase.

\textbf{Methodological point:} Quantum Graphity treats locality as
emergent. FFGFT treats it as given. This is not a contradiction but
a difference of description level -- in FFGFT the underlying
geometry is already chosen; in Quantum Graphity the choice of
geometry is left to the system.

% ==============================================================================
\section{Reduction 3: windings to spin labels}
% ==============================================================================

\subsection{FFGFT: specific winding quantum numbers}

In FFGFT the winding quantum numbers
$(n_x, n_y, n_z, k_t) \in \mathbb{Z}^4$ are \emph{physically
interpreted}:
\begin{itemize}[leftmargin=2em,itemsep=2pt,topsep=2pt]
  \item $(n_x, n_y, n_z)$ encode spatial mode position and angular
        momentum (Doc.~210),
  \item $k_t$ carries the mass energy of the mode
        ($m c^2 = \hbar f(k_t)$),
  \item specific tuples yield specific Standard-Model particles.
\end{itemize}

\subsection{\texorpdfstring{Quantum Graphity: fermionic occupation and $(j, m)$ labels}{Quantum Graphity: fermionic occupation and (j,m) labels}}

In Quantum Graphity the edges in the \emph{base model} carry only a
fermionic occupation number $n_{ab} \in \{0, 1\}$ -- off or on. In
the \emph{extended version} (§IV of the paper) the \enquote{on}
states acquire additional labels $|j, m\rangle$ with $j = 0, 1$ and
$m \in \{-j, \ldots, +j\}$. The labels are general
\emph{spin-network-like} and belong to the LQG tradition. The
operators $M, M^\pm$ obey the usual angular-momentum algebra
\[
  [M^+, M^-] = M, \qquad [M, M^\pm] = \pm M^\pm.
\]

\subsection{Reduction}

The reduction FFGFT $\to$ Quantum Graphity in this layer is
\emph{forgetting the specific physical meaning} of the winding
quantum numbers. The specific tuple $(n_x, n_y, n_z, k_t)$ that for
example describes an electron is replaced by:
\begin{itemize}[leftmargin=2em,itemsep=2pt,topsep=2pt]
  \item In the base model: only \enquote{on}/\enquote{off} per edge --
        no quantum-number information at all.
  \item In the extended version: $(j, m)$ labels carrying general
        spin-1 information, without specifying which number carries
        which physical quantity.
\end{itemize}

\textbf{What is lost:} the connection to the Standard Model. In
FFGFT the winding structure directly yields the mass formula
$m_i = r_i \xi^{p_i} v$ and thus all 12 fermion masses. In Quantum
Graphity there is no mechanism that links specific $(j, m)$ values
to specific masses.

% ==============================================================================
\section{Reduction 4: geometry parameter to phase parameters}
% ==============================================================================

\subsection{\texorpdfstring{FFGFT: $\xi_0 = 4/30000$ is geometrically fixed}{FFGFT: xi0 = 4/30000 is geometrically fixed}}

In FFGFT the geometry constant $\xi_0 = 4/30000$ is not a free
parameter, but follows from Higgs matching
$\xi_0 = \lambda_h^2 v^2/(16\pi^3 m_h^2)$ (Doc.~006). It is
geometrically and physically \emph{uniquely fixed}.

\subsection{Quantum Graphity: multiple coupling parameters}

In Quantum Graphity there is a whole set of coupling parameters:
$g_V, g_B, g_C, g_D, g_\pm$ (Hamiltonian components) plus $v_0$
(preferred vertex valence) and $p, r$. Their values determine where
the phase transition occurs, which subgraph emerges in the
low-temperature phase (typically $v_0 = 3$ yields a \emph{hexagonal}
configuration), and how strong the string-net condensation is.
None of these parameters is fixed externally -- they are \emph{free}.

\subsection{Reduction}

The transition FFGFT $\to$ Quantum Graphity in this layer is the
\emph{multiplication and generalisation of parameters}. The single
geometrically fixed quantity $\xi_0 = 4/30000$ is replaced by a
whole set of free couplings.

\textbf{What is lost:} predictive power. The values of all SM
constants in FFGFT follow from the single parameter $\xi_0$. If
one replaces $\xi_0$ by several free couplings, these predictions
are lost -- one can determine the phase diagram, but not the values
of physical constants.

\textbf{Methodological point:} Quantum Graphity shows that a family
of parameters suffices to generate a rich emergent geometry. FFGFT
shows additionally that all these parameters can be fixed by a
single geometric quantity ($\xi_0$) -- and then also yield all SM
constants.

% ==============================================================================
\section{Reduction 5: dropping the matter sector}
% ==============================================================================

\subsection{FFGFT: full matter sector}

In FFGFT, all 12 fermion masses, the fine-structure constant
$\alpha = \xi \cdot E_0^2$, the Higgs self-coupling, and the
spin-orbit constants of atomic physics follow from the winding
quantum numbers (Doc.~006, 174, 210). The matter sector is not
introduced separately, but \emph{integral part} of the geometry.

\subsection{Quantum Graphity: no matter sector (in the base model)}

In the base model of Quantum Graphity there is no matter
specification -- only the binary edge states $|0\rangle, |1\rangle$
and the emergent lattice. In the \emph{extended version} with
$(j, m)$ labels a \emph{$U(1)$ gauge field} emerges (through the
string-net condensation of Levin-Wen) -- this is analogous to a
photon field. But there are no fermions, no quarks, no specific
masses -- the string-net condensation does generate point-like
excitations with mass $\propto g_C$, but these masses are free
parameters, not determined by a geometric source.

\subsection{Reduction}

The final reduction step is the complete \emph{discarding} of the
matter-determining FFGFT sector. FFGFT contains matter as winding
modes on $T^4$ with \emph{fixed mass relations}. Quantum Graphity
in the extended version contains emergent $U(1)$ gauge fields and
weakly bound excitations, but without a geometric mass formula.

\textbf{What is lost:} everything FFGFT delivers for the explanation
of SM masses and couplings.

\textbf{What remains:} the emergent gauge theory. The fact that
$U(1)$ emerges in both theories from microstructure is a real
structural agreement -- it is the only row in the translation
table that truly \emph{exactly} corresponds.

% ==============================================================================
\section{Translation table}
% ==============================================================================

\begin{center}
\renewcommand{\arraystretch}{1.4}
{\small
\adjustbox{max width=\textwidth}{%
\begin{tabular}{>{\raggedright\arraybackslash}p{3.1cm}>{\raggedright\arraybackslash}p{3.7cm}>{\raggedright\arraybackslash}p{3.3cm}}
\toprule
\textbf{FFGFT structure} & \textbf{Quantum Graphity counterpart} & \textbf{Reduction} \\
\midrule
4D-torus $T^4$ & complete graph $K_N$ & discretisation + topology forgetting \\[3pt]
windings $(n_x, n_y, n_z, k_t)$ & occupation $n_{ab}$ + ext.\ $(j, m)$ & specific meaning lost \\[3pt]
mode functions $\varphi_{n,l,j}$ & edge states $|0\rangle, |1\rangle$ & smooth $\to$ discrete \\[3pt]
$\xi_0 = 4/30000$ single constant & $g_V, g_B, g_C, g_D, g_\pm$ and $v_0$ free & predictive power lost \\[3pt]
low-energy smoothing to $T^4$ geometry & low-temperature: hexagonal $3$-regular & same principle \\[3pt]
emergent $U(1)$ from EM winding & emergent $U(1)$ via Levin-Wen (ext.\ only) & true correspondence \\[3pt]
12 fermion masses & --- & no counterpart \\[3pt]
$\alpha = \xi E_0^2$ & --- & no counterpart \\[3pt]
spin-orbit constants & --- & no counterpart \\
\bottomrule
\end{tabular}%
}
}
\renewcommand{\arraystretch}{1.0}
\end{center}

\textbf{Status of the correspondences:}
\begin{itemize}[leftmargin=2em,itemsep=2pt,topsep=2pt]
  \item \textbf{One real correspondence:} emergent $U(1)$ in both
        theories from microstructure (in QG only in the extended
        version with $(j, m)$ labels).
  \item \textbf{Structural parallels:} low-energy smoothing to
        ordered geometry is the same basic principle in both
        theories.
  \item \textbf{Reduction losses:} all SM-specific quantities
        (masses, $\alpha$, spin-orbit) have no Quantum-Graphity
        counterpart -- they are lost in the reduction process.
\end{itemize}

% ==============================================================================
\section{What is lost in reduction -- and what remains}
% ==============================================================================

\subsection{What is lost}

The five reduction steps lead to a progressive loss:
\begin{itemize}[leftmargin=2em,itemsep=2pt,topsep=2pt]
  \item differential geometry (Reduction 1),
  \item locality as a given structure (Reduction 2),
  \item physical meaning of the winding quantum numbers
        (Reduction 3),
  \item geometry-parameter specificity (Reduction 4),
  \item entire matter sector (Reduction 5).
\end{itemize}

\subsection{What remains}

Remarkably, \emph{the essential geometry-emergence scheme} is
preserved:
\begin{itemize}[leftmargin=2em,itemsep=2pt,topsep=2pt]
  \item In FFGFT: the full geometry is not directly visible but
        emerges from the mode configuration in the low-energy limit.
  \item In Quantum Graphity: geometry is not laid down in advance
        but emerges from the phase transition in the low-temperature
        phase.
\end{itemize}

\textbf{Both theories share the same methodological position:
geometry is not given but emergent.} FFGFT specifies \emph{which}
geometry ($T^4$) and \emph{why} (geometric Higgs matching,
$\xi_0$ quantisation). Quantum Graphity shows that this idea of an
emergent geometry from discrete microstructure is \emph{not exotic},
but is pursued in established peer-reviewed quantum-gravity research.

% ==============================================================================
\section{Methodological position}
% ==============================================================================

\subsection{What this reduction shows}

The five reduction steps show that Quantum Graphity is a
\emph{subset of FFGFT} -- analogous to the Bloch sphere, but
through different reduction steps. Concretely:
\begin{itemize}[leftmargin=2em,itemsep=2pt,topsep=2pt]
  \item Bloch sphere $\subset$ FFGFT through \emph{geometric
        reduction} (pole contraction, loop cutting).
  \item Quantum Graphity $\subset$ FFGFT through \emph{structural
        reduction} (discretisation, topology forgetting, matter
        dropping).
\end{itemize}

\subsection{Why this matters}

This strategy -- a \emph{subset relation} instead of a full
bijection -- is methodologically important:

\begin{itemize}[leftmargin=2em,itemsep=2pt,topsep=2pt]
  \item It is \textbf{honest:} no claim of an exact translation
        where only a partial one is possible.
  \item It \textbf{positions FFGFT as the more comprehensive
        framework:} Quantum Graphity is a special case, not a
        competitor.
  \item It \textbf{connects FFGFT to established research:} FFGFT
        shares the methodological core programme with Konopka,
        Markopoulou, Severini -- and thus also with Loop Quantum
        Gravity, Causal Dynamical Triangulation, Wolfram
        hypergraphs.
  \item It \textbf{shows FFGFT's added value:} while Quantum
        Graphity delivers only emergent geometry, FFGFT delivers
        additionally the Standard-Model constants.
\end{itemize}

\subsection{Connection to Doc.~230 and 231}

\begin{itemize}[leftmargin=2em,itemsep=2pt,topsep=2pt]
  \item \textbf{Doc.~230} showed: Hilbert-space QM $\subset$ FFGFT
        (through geometric reduction to the Bloch sphere). This
        reduction is \emph{carried out} -- the translation bridge
        $\alpha = \sqrt{(1+z)/2}\,e^{i\theta/2}$ is a concrete
        bijection.
  \item \textbf{Doc.~231} showed: Hilbert-space QM can be
        \emph{extended} to full FFGFT (through four mathematical
        structures: $SU_q(2)$, winding quantum numbers, spinor
        bundles, Kaluza-Klein). These extensions are
        \emph{established} -- $SU_q(2)$ exists since 1985, spinor
        bundles since the 1950s, Kaluza-Klein since 1921.
  \item \textbf{Doc.~232 (this document)} sketches: Quantum Graphity
        $\subset$ FFGFT (through structural reduction). This
        reduction is \emph{not yet carried out} -- it is plausibly
        sketched, but not rigorously constructed (see §10.4).
\end{itemize}

Together the three documents yield the picture: \textbf{FFGFT
demonstrably contains Hilbert-space QM as a subset (Doc.~230/231)
and plausibly contains the discrete-geometry programs as special
cases, provided the triangulation reduction can actually be carried
out (Doc.~232).} It is not a competitor to these theories, but the
relationship to Hilbert-space QM and to Quantum Graphity is
\emph{not symmetric} -- see §10.4 and §11.

\subsection{Status of the reduction: hypothesis, not proof}

\textbf{Important caveat.} At several points in the FFGFT documents
-- including this one -- there is talk of \enquote{triangulations},
\enquote{tetrahedron geometry}, and \enquote{translatability into
triangle or network descriptions}. This language is useful for
locating FFGFT methodologically within the broader quantum-gravity
tradition (LQG, CDT, Wolfram, Quantum Graphity), but it must not
be taken as a proved mathematical translation theorem.

\textbf{What actually holds:}
\begin{itemize}[leftmargin=2em,itemsep=2pt,topsep=2pt]
  \item \emph{Methodological kinship:} FFGFT shares with Quantum
        Graphity (and LQG, CDT) the basic programme \enquote{emergent
        geometry from discrete microstructure}. This is a conceptual
        statement, not a mathematical theorem.
  \item \emph{Plausible reduction chain:} the five reduction steps
        in §3-7 of this document are a plausible sketch of how a
        translation FFGFT $\to$ Quantum Graphity \emph{could look}.
        It is consistent with standard differential topology
        (simplicial approximation) and with the structures of the
        Quantum-Graphity Hilbert space (fermionic tensor product
        over edges).
  \item \emph{Internal FFGFT tetrahedral geometry:} The tetrahedron
        factor $4/3$ in $\xi_0 = 4/30000$ is an \emph{internal}
        FFGFT statement (Doc.~180, 202), not a claim about
        translatability to LQG spin-foam models.
\end{itemize}

\textbf{What remains open:}
\begin{itemize}[leftmargin=2em,itemsep=2pt,topsep=2pt]
  \item A \emph{rigorous convergence proof} that a specific $T^4$
        triangulation with $N \to \infty$ approximates the FFGFT
        mode functions $\varphi_{n,l,j}$.
  \item A \emph{Hamiltonian translation theorem} relating the FFGFT
        operator $\mathcal{D}_{T^4}$ to the Quantum-Graphity
        Hamiltonian chain $H_V + H_B + \cdots$.
  \item A \emph{recovery theorem} showing that the low-temperature
        phase of a Quantum-Graphity model with appropriately chosen
        couplings yields an effective $T^4$.
\end{itemize}

\textbf{Methodological consequence:} The statement \enquote{Quantum
Graphity $\subset$ FFGFT} is a \emph{plausible programme item}, not
a proved structural statement -- in contrast to the Hilbert-space
translation in Doc.~230, which exists as a concrete bijection.
This asymmetry between Docs 230 and 232 must be kept in mind when
both documents are used. Anyone presenting FFGFT's relationship to
LQG, CDT, Wolfram or Quantum-Graphity programmes should formulate
it as \emph{methodological kinship with shared basic assumptions} --
not as a formally proved reduction.

% ==============================================================================
\section{Comparison of explanatory power: Hilbert-space QM vs.\ Quantum Graphity}
% ==============================================================================

The reduction FFGFT $\to$ Quantum Graphity differs \emph{fundamentally}
from the reduction FFGFT $\to$ Hilbert-space QM (Doc.~230) -- not only
in the reduction steps, but in the \textbf{explanatory power} of the
respective target theory. This asymmetry is methodologically
important and is made explicit here.

\subsection{Hilbert-space QM: substantial content}

Standard Hilbert-space quantum mechanics is a \emph{mathematically
mature, complete} theory:
\begin{itemize}[leftmargin=2em,itemsep=2pt,topsep=2pt]
  \item It describes matter completely (wave functions, spinors,
        many-particle states).
  \item It yields precise predictions (atomic spectra, scattering
        amplitudes, quantum gate operations).
  \item It is overwhelmingly experimentally confirmed.
  \item It accommodates all three gauge groups
        $SU(3) \times SU(2) \times U(1)$ in tensor-product form.
  \item It has a century of mathematical development behind it.
\end{itemize}

The translatability FFGFT $\leftrightarrow$ Hilbert-space QM
(Doc.~230) is therefore a \emph{bridge between two mature theories},
with small, quantifiable loss ($K_{\text{frak}}$ correction,
$\Delta\text{CHSH} \approx 10^{-5}$).

\subsection{Quantum Graphity: programmatic content}

Quantum Graphity is a \emph{programme}, not a complete tool:
\begin{itemize}[leftmargin=2em,itemsep=2pt,topsep=2pt]
  \item Matter is not described in the base model; in the extended
        version only as $(j, m)$ labels without mass specificity.
  \item Predictions are \emph{qualitative} (phase transition,
        emergent locality), not quantitative.
  \item There are no direct experimental confirmations -- Quantum
        Graphity operates at the Planck scale.
  \item Only one gauge group ($U(1)$) emerges, and only in the
        extended version with Levin-Wen string-net condensation.
  \item It is 17~years old (first paper 2008), with a small research
        community.
\end{itemize}

\subsection{Quantitative comparison}

\begin{center}
\renewcommand{\arraystretch}{1.4}
{\small
\adjustbox{max width=\textwidth}{%
\begin{tabular}{>{\raggedright\arraybackslash}p{3.4cm}>{\raggedright\arraybackslash}p{3.4cm}>{\raggedright\arraybackslash}p{3.4cm}}
\toprule
\textbf{Aspect} & \textbf{Hilbert-space QM} & \textbf{Quantum Graphity} \\
\midrule
Matter sector & fully described & not present (base model) \\[3pt]
Masses & external Yukawa parameters & not present \\[3pt]
Gauge theories & $SU(3)\!\times\!SU(2)\!\times\!U(1)$ & only $U(1)$ (ext.\ version) \\[3pt]
Spin algebra & complete & added if needed \\[3pt]
Predictions & precise quantitative & qualitative \\[3pt]
Experimental confirmation & overwhelming & none direct \\[3pt]
Mathematical maturity & one century & 17~years \\[3pt]
Reduction loss vs.\ FFGFT & small, quantifiable & massive \\
\bottomrule
\end{tabular}%
}
}
\renewcommand{\arraystretch}{1.0}
\end{center}

\subsection{What the asymmetry means for FFGFT's position}

The two reductions show FFGFT's relationship to two very different
theory classes:

\begin{itemize}[leftmargin=2em,itemsep=2pt,topsep=2pt]
  \item \textbf{Towards Hilbert-space QM:} FFGFT has to respect a
        \emph{large, established} theory -- and does so, with full
        translatability and small geometric corrections. This is
        FFGFT's proof of \emph{consistency} with known physics.

  \item \textbf{Towards Quantum Graphity:} FFGFT contains a
        \emph{small, programmatic} model as a special case. This is
        FFGFT's proof of \emph{methodological kinship} with other
        quantum-gravity programmes -- without FFGFT having to inherit
        their limitations.
\end{itemize}

\textbf{The point:} Quantum Graphity is \emph{not wrong}, but
\emph{little}. FFGFT shares with it the basic methodological
programme \enquote{emergent geometry from discrete microstructure},
but goes far beyond: it specifies the geometry ($T^4$), the single
parameter ($\xi_0$), and from these yields all 12 fermion masses,
the fine-structure constant, the gauge-group structure, and the
spin-orbit constants -- all parameter-free.

\subsection{Strategic consequence}

\begin{itemize}[leftmargin=2em,itemsep=2pt,topsep=2pt]
  \item \textbf{To Hilbert-space practitioners (mathematicians,
        QM theorists):} FFGFT respects full standard QM and extends
        it consistently (Doc.~230, 231). The translation tables are
        interoperable.

  \item \textbf{To quantum-gravity researchers (LQG, Wolfram,
        Quantum Graphity):} FFGFT shares the basic programme
        \enquote{emergent geometry from microstructure} but goes
        further and yields SM constants parameter-free. Quantum
        Graphity is a special case in which almost all FFGFT
        statements are dropped.
\end{itemize}

\textbf{FFGFT's strength arises precisely from this asymmetry:}
it is \emph{large enough} to contain Hilbert-space QM, and
\emph{specific enough} to extend Quantum Graphity.

% ==============================================================================
\section{Summary}
% ==============================================================================

\textbf{Quantum Graphity is plausibly representable as a subset of
FFGFT -- in the sense of a sketched, not carried-out reduction.}

The reduction chain FFGFT $\to$ Quantum Graphity consists of five
steps:
\begin{enumerate}[leftmargin=2em,itemsep=2pt,topsep=2pt]
  \item Discretisation of the $T^4$ continuum (sketched as a
        triangulation -- not yet rigorously constructed).
  \item Forgetting the $T^4$ topology in favour of the complete
        graph $K_N$.
  \item Generalisation of the winding quantum numbers
        $(n_x, n_y, n_z, k_t)$ to general edge states ($n_{ab}$ in
        the base model, $(j, m)$ in the extended version).
  \item Generalisation of the geometry parameter
        $\xi_0 = 4/30000$ to a set of free couplings $g_V, g_B,
        g_C, g_D, g_\pm, v_0, \ldots$.
  \item Dropping the matter sector (masses, $\alpha$, spin-orbit).
\end{enumerate}

What remains after reduction: the \emph{geometry-emergence scheme}
and the \emph{emergent $U(1)$ gauge theory}. What is lost: the
SM-specific predictions.

\textbf{Methodological position:} FFGFT shares with Quantum Graphity
(and LQG, CDT, Wolfram hypergraphs) the basic programme
\enquote{emergent geometry from discrete microstructure}. FFGFT
goes further: it specifies the geometry ($T^4$), the parameter
($\xi_0$), and from that yields all SM constants parameter-free.

\textbf{Status of the reduction:} In contrast to the Hilbert-space
translation in Doc.~230 (concrete bijection with
$\alpha = \sqrt{(1+z)/2}\,e^{i\theta/2}$), the reduction in this
document is a \emph{hypothesis} -- plausibly supported by standard
differential topology, but without a carried-out convergence proof.
The statement \enquote{Quantum Graphity $\subset$ FFGFT} is therefore
to be understood as a \emph{methodological programme item}, not as
a formally proved structural statement.

Quantum Graphity is therefore \emph{not a contradiction} of FFGFT
but a \emph{partial FFGFT skeleton} -- the emergent-geometry
component without the matter specificity. The existence of Quantum
Graphity as a peer-reviewed research programme demonstrates that
FFGFT's methodological assumptions are \emph{not exotic}.
